% !TEX root = ../main.tex
\chapter{Related Work}
\label{chap:related_work}

Digital DNA storage has progressed through several coding architectures since the first large-scale demonstrations in 2012. This chapter traces that evolution, reviews the theoretical capacity bounds that frame the field, surveys the existing software tools and positions the simulator developed here relative to prior work.

\section{Coding Schemes for DNA Data Storage}
\label{sec:related_dna_encoding}

This section reviews the primary coding schemes in chronological order, highlighting key design choices, practical constraints and reported storage densities.

\subsection{Church, Gao and Kosuri (2012)}

The work by \citeonline{church2012} is considered one of the initial milestones in scaled digital DNA storage. The authors encoded an entire 53,426-word book, including images and JavaScript code, into synthetic oligonucleotides. The encoding scheme employed was straightforward: each bit of the original file was directly mapped to a nitrogenous base (0 $\rightarrow$ A or C, 1 $\rightarrow$ G or T), alternating assignments to avoid long homopolymers.

However, the approach of \citeonline{church2012} did not employ any mathematical error-correcting redundancy beyond physical replicas of the oligonucleotides. Redundancy was introduced exclusively by synthesizing multiple copies of each sequence, increasing synthesis costs. In addition, the achieved storage density was only 0.83 bits per nucleotide, well below the theoretical unconstrained limit of $\log_2 4 = 2.0$ bits per nucleotide. This limitation stemmed from the 1-bit-per-base mapping, which discarded half of the informational capacity of DNA's quaternary alphabet.

From a fault-tolerance perspective, the scheme exhibited significant fragility, as the loss of a single oligonucleotide resulted in the irreversible loss of all bits contained within it, with no recovery mechanism. Despite these limitations, the work demonstrated the practical feasibility of the concept and motivated numerous subsequent advances.

\subsection{Goldman et al. (2013)}

The following year, \citeonline{goldman2013} proposed a more sophisticated encoding scheme published in \emph{Nature}. The authors utilized a base-3 (trit) mapping with alphabet rotation, where each trit was mapped to one of three bases distinct from the preceding base, preventing consecutive base repetitions by construction. The rotation scheme reduced the occurrence of homopolymers by construction, preventing runs of length greater than 1.

\citeonline{goldman2013} also introduced redundancy through overlapping segments: each region of the file was covered by four partially overlapping oligonucleotides, enabling data recovery even in the presence of individual segment loss. Although the theoretical encoding potential of the rotating scheme was 1.58 bits per nucleotide (disregarding redundancy), the effective net density resulting from four-fold overlapping and indexing overhead in that experimental configuration was approximately $0.33$ bits per nucleotide, which does not constitute an invariant upper bound for all overlap-based schemes. The four-fold overlap improved tolerance to segment losses, although the resulting net density in that trial remained well below channel capacity.

At the same time, overlap redundancy operated at a fixed rate set at design time, limiting adaptability to channels with varying loss profiles. If physical degradation exceeded the predefined redundancy budget, decoding failed irreversibly without the capability to request or generate additional segments.

\subsection{Grass et al. (2015)}

\citeonline{grass2015} approached long-term storage from a distinct perspective, focusing on physical DNA preservation via silica encapsulation. From an information-theoretic standpoint, the authors employed Reed--Solomon codes over the finite field $GF(47)$ (mapping field elements to base triplets to avoid homopolymer runs), organized into two hierarchical tiers: an inner code to correct substitution errors within each oligonucleotide and an outer code to recover lost oligonucleotides.

This hierarchical strategy treated intra-oligo substitutions and whole-sequence erasures as distinct error processes. The authors experimentally demonstrated error-free recovery of 83 kilobytes of data after accelerated aging equivalent to 2,000 years under ambient conditions.

Nonetheless, like \citeonline{goldman2013}, the scheme of \citeonline{grass2015} relied on fixed-rate coding: redundancy was predetermined statically and could not adapt dynamically to channel conditions. In addition, general Reed--Solomon decoding has quadratic complexity, posing computational bottlenecks for larger data volumes.

\subsection{Bornholt et al. (2016) and Blawat et al. (2016)}

Two works published in 2016 advanced in complementary directions. \citeonline{bornholt2016} proposed a complete end-to-end DNA storage architecture, introducing random-access capabilities via selective PCR primers. The system demonstrated the ability to retrieve specific data subsets without sequencing the entire oligonucleotide pool, bringing functionality closer to conventional random-access storage systems.

In parallel, \citeonline{blawat2016} systematically investigated Forward Error Correction (FEC) codes tailored to the DNA channel, employing adapted Reed--Solomon codes. The authors evaluated tradeoffs between code rate and correction capability, demonstrating that carefully tuning error-correction parameters was decisive for overall storage efficiency.

Both works, however, used fixed-rate coding whose redundancy could not adapt to varying channel conditions.

\subsection{Erlich and Zielinski: DNA Fountain (2017)}

The influential work of \citeonline{erlich2017}, published in \emph{Science} in 2017, pioneered the application of \emph{fountain codes} to DNA data storage. The DNA Fountain architecture combines a two-tier hierarchical coding structure: an outer rateless Luby Transform (LT) code operating across blocks to reconstruct the file in the presence of full oligonucleotide erasures via the Robust Soliton Distribution (RSD), and an inner Reed--Solomon code (2 parity bytes per oligonucleotide) designed to detect and discard localized substitution and indel errors caused by synthesis and sequencing. The rateless property of the outer code enables generating encoded droplets on demand until the receiver collects a sufficient quantity ($K(1 + \epsilon)$) to reconstruct the original data via the iterative \emph{peeling decoder}.

DNA Fountain introduced three key advantages over earlier schemes:

\begin{itemize}
    \item \textbf{Rateless capability.} Unlike fixed-rate codes (such as Reed--Solomon or overlapping windows), DNA Fountain generates droplets continuously until the desired number of valid oligonucleotides is gathered. Rateless encoding enables generating additional oligonucleotides digitally and tuning the redundancy budget without requiring real-time feedback retransmission. However, an already synthesized physical pool remains limited to the generated batch.
    
    \item \textbf{Reported storage density.} At the time of publication, the reported net density of 1.57 bits per nucleotide was among the highest values in the DNA storage literature, representing approximately 86\% of the practical Shannon capacity estimated under unordered pooling and indexing constraints ($\approx 1.83$ bits/nt), compared to a theoretical coding capacity of 1.98 bits/nt under GC and homopolymer constraints and the unconstrained limit of $\log_2 4 = 2.0$ bits per nucleotide.
    
    \item \textbf{Integrated biological screening.} The process of discarding and re-drawing candidate oligonucleotides that violate biochemical synthesis constraints (such as GC bias or long homopolymers) is naturally absorbed by the rateless nature of fountain codes, without structural penalties to the coding scheme.
\end{itemize}

The authors experimentally validated error-free encoding and decoding of a 2.14-megabyte file (including a full operating system, a computer virus and a \$50 Amazon gift card) partitioned into $K = 67{,}088$ source segments, packaged into 72,000 oligonucleotides of 200 nucleotides each ($\approx 7.3\%$ overhead).

\subsection{Organick et al. (2018)}

\citeonline{organick2018} expanded the scale of DNA data storage to hundreds of megabytes, encoding 200 megabytes across more than 13 million oligonucleotides. The work demonstrated large-scale random access using selective primer libraries, enabling the retrieval of individual files from a complex heterogeneous pool without sequencing the entire molecular library.

Regarding coding strategy, \citeonline{organick2018} employed a hierarchical approach based on concatenated Reed--Solomon codes (an inner code per oligo and an outer code across oligos) with fixed redundancy, distinct from the rateless methodology of DNA Fountain. While demonstrating scalability to hundreds of megabytes, error-free recovery depended heavily on high sequencing coverage, and redundancy overhead remained statically defined at design time.

\section{Theoretical Foundations and Capacity Bounds}
\label{sec:related_theoretical}

In parallel with experimental breakthroughs, theoretical studies formalized fundamental limits for DNA data storage, establishing tradeoffs among capacity, redundancy and reliability.

\citeonline{heckel2017} derived fundamental capacity limits for the DNA storage channel under a model where sequences are stored unordered in a molecular pool and subjected to random sampling and mutation errors. The authors proved that, for capacity to remain positive, oligonucleotide length must scale at least logarithmically with the number of stored molecules, mathematically formalizing the penalty imposed by unordered pooling. This result provides motivating theoretical context for unordered molecular channels, rather than physical validation of the simplified memoryless Binary Erasure Channel (BEC) adopted in this work's simulations.

More recently, \citeonline{shomorony2022} published a review of DNA data storage foundations in \emph{Foundations and Trends in Communications and Information Theory}, consolidating information-theoretic principles applicable to the domain. That work formalized the distinction between the \emph{sequencing channel} (where ordering information is lost) and the \emph{storage channel} (involving physical molecular errors), establishing asymptotic capacity bounds and the conditions under which indexing and erasure coding approach channel capacity in abstract molecular storage models.

These theoretical erasure channel models provide reference bounds for evaluating the relationship between overhead and decoding success rate observed in the experimental results of Chapter~\ref{chap:results}.

\section{Software Tools and Simulators}
\label{sec:related_simulators}

Despite growing interest in DNA data storage, open and modular software tools designed for academic experimentation remain limited in the literature. The tools identified in the reviewed literature are summarized below.

In rateless coding literature, degree distribution optimization for short message regimes was investigated by \citeonline{hyytia2007}. The authors analyzed LT codes for small block counts in telecommunication erasure channels, analytically and heuristically demonstrating that Luby's standard RSD suffers efficiency degradation in finite regimes, and proposed alternative distributions via absorbing Markov chains and recursive combinatorial optimization. That study focused on classical network channels, without addressing biochemical sequence constraints or DNA screening dynamics.

The storage architecture proposed by \citeonline{bornholt2016} included a software pipeline for encoding and decoding geared toward end-to-end wet-lab validation (synthesis, storage, sequencing), with software closely coupled to their specific experimental setup.

\citeonline{press2020} developed the HEDGES (\emph{Hash Encoded, Decoded by Greedy Exhaustive Search}) system, an error-correcting code designed specifically for the DNA channel capable of correcting insertions and deletions in addition to substitutions. The HEDGES implementation represented an important contribution, focusing primarily on intra-strand indels and localized errors.

\citeonline{chandak2019} explored Low-Density Parity-Check (LDPC) codes for DNA storage, evaluating tradeoffs between synthesis costs and sequencing coverage, representing a structural approach distinct from rateless fountain codes.

\citeonline{antkowiak2020} proposed a low-cost photolithographic synthesis pipeline coupled with dedicated alignment and reconstruction algorithms tailored to their proprietary physical setup.

\citeonline{erlich2017} released the source code of their implementation in a public repository under the GPL-3.0 license. The original codebase was built as a pipeline comprising Python 2 modules and compiled Cython routines for fast execution (\texttt{fountain.pyx}, \texttt{glass.pyx}, \texttt{robust\_solition.pyx}, \texttt{lfsr.py}, \texttt{encode.py} and \texttt{receiver.py}). The tool exposes a command-line interface with parameters such as block size, GC boundaries (default range 45\% to 55\%), failure probability $\delta$, tuning parameter $c$ and oligonucleotide target counts. However, this codebase was designed as a companion artifact to facilitate their large-scale wet-lab experiment (packaging multi-megabyte archives with $K = 67{,}088$ source blocks, 2-byte Reed--Solomon parity per oligo, and 32-bit headers). The original repository focuses on reproducing that large-scale experimental pipeline; it does not provide the specific decoding instrumentation of ripple dynamics, standard-library execution, or automated batch parametric grid execution with CSV metric logging implemented here.

\citeonline{schwarz2021} introduced \textbf{NOREC4DNA}, an open-source Python framework designed for comparative analysis of rateless erasure codes for DNA data storage. NOREC4DNA implements LT, Raptor and Online codes, supporting ideal and robust soliton distributions as well as programmable adaptive distributions. The framework incorporates biological constraints and noisy channel simulation with substitution and erasure errors. More recently, \citeonline{schwarz2024} advanced numerical optimization of rateless codes for DNA storage, presenting algorithms to tailor degree distributions to specific molecular error profiles and sequence constraints, achieving measurable capacity and reliability gains over generic distributions.

\section{Identified Gap and Positioning of this Work}
\label{sec:identified_gap}

A review of the works surveyed above shows that the reviewed approaches emphasize different goals: large-scale wet-lab replication artifacts (such as the original DNA Fountain~\cite{erlich2017}), theoretical analyses of short codes in telecommunication channels~\cite{hyytia2007}, and broad comparative frameworks focused on advanced rateless codes and numerical distribution optimizations~\cite{schwarz2021, schwarz2024}.

To address these limitations, the simulator developed in this work focuses on four core priorities:

\begin{enumerate}
    \item \textbf{Modular and self-contained implementation.} An object-oriented Python~3 implementation relying exclusively on the standard library, eliminating compiled C or Cython dependencies, and facilitating straightforward replication in educational and research environments.
    
    \item \textbf{Analytical peeling instrumentation.} Direct exposure of the internal mechanics of the peeling decoder, logging step-by-step \emph{ripple} queue trajectories, exact XOR operation counts per source block and droplet degree transitions on the bipartite graph.
    
    \item \textbf{Experimental evaluation in short-block regimes ($K \le 64$).} Empirical characterization of the combined behavior of the classical RSD under biochemical screening constraints and stochastic erasure channels across a factorial grid of 216 configurations, mapping $\tau(d)$ mode transitions and ripple exhaustion conditions.
    
    \item \textbf{Batch experimental infrastructure.} Automated execution scripts for multidimensional parameter grids with pseudorandom seed control and direct CSV metrics export.
\end{enumerate}

Table~\ref{tab:comparacao_especificacao} compares the technical specifications of the original DNA Fountain by \citeonline{erlich2017} and the NOREC4DNA framework by \citeonline{schwarz2021} with the simulator developed in this work.

\begin{table}[htbp]
    \centering
    \small
    \caption{Technical comparison among the original DNA Fountain, NOREC4DNA and the simulator developed in this work.}
    \label{tab:comparacao_especificacao}
    \resizebox{\textwidth}{!}{%
    \begin{tabular}{lp{4.0cm}p{4.5cm}p{4.2cm}}
        \toprule
        \textbf{Feature} & \textbf{Original DNA Fountain}~\cite{erlich2017} & \textbf{NOREC4DNA}~\cite{schwarz2021} & \textbf{This Work's Simulator} \\
        \midrule
        Language & Python 2.7 and Cython (C extensions) & Python 3 (requires NumPy, SciPy) & Pure Python 3.12 (standard library only) \\
        License & GPL-3.0 & Open Source (AGPL-3.0) & Open Source (MIT License) \\
        Inner code (per oligo) & Reed--Solomon (2 bytes/oligo) & Configurable (none, CRC, or RS) & CRC-8 corruption detection (non-correcting) \\
        Outer code (across blocks) & LT Code (RSD) & LT, Raptor and Online codes & LT Code (RSD) and $GF(2)$ Gaussian validator \\
        Seed header size & 32 bits (4 bytes) in header & Configurable (defaults to 16 bits) & 16 bits (2 bytes) with XOR mask \\
        Target regime ($K$) & Large scale ($K = 67{,}088$) & Broad range (small to large) & Short blocks ($K \in \{16, 32, 64\}$) \\
        GC content range & 45\% to 55\% & Configurable biochemical filter & 40\% to 60\% ($H_{\max} \le 3$) \\
        Payload handling & LFSR scrambler and PCR primers & Scrambler and adaptable mapping & Additive scrambler (XOR mask) and direct mapping \\
        Core purpose & Software pipeline used in a wet-lab demonstration & Comparative benchmarking framework & Analytical and didactic instrumentation \\
        Collected metrics & File reconstruction success & Throughput, code rate, recovery & XOR operations, ripple queue, runtimes, erasures \\
        \bottomrule
    \end{tabular}%
    }
    \fonte{Prepared by the author based on \citeonline{erlich2017} and \citeonline{schwarz2021} (2026).}
\end{table}
